% ============================================================
%  Chapter 5: The Partition Lagrangian
%  Source: bounded-phase-space-trajectory.tex, Consequence 15
% ============================================================

\epigraph{%
  The equation of motion is not a force law.\\
  It is gradient descent on the partition depth.%
}{}

\section{Deriving the Lagrangian from the Axiom}
\label{sec:ch5-derive}

Classical mechanics is usually \emph{defined} by a Lagrangian or Hamiltonian,
taken as a given for the physical system in question.
Here we derive the unique Lagrangian consistent with BPSL from first
principles, without assuming Newton's second law.

The argument has three steps:
\begin{enumerate}
  \item Identify the scalar field that the system naturally extremises.
  \item Write the most general Lagrangian coupling a particle to that field
    that is consistent with Galilean invariance and completeness of
    scalar-plus-vector coupling.
  \item Show that the Euler--Lagrange equations take the Lorentz-force shape
    but are force-free.
\end{enumerate}

\section{The Partition Depth Field}
\label{sec:ch5-depth}

The partition algebra $\mathcal{A}_\infty$ assigns to every point $\mathbf{x}$
in configuration space a partition depth --- the finest partition level at
which $\mathbf{x}$ is an interior point of a cell.
We define the \emph{partition depth field}:
\begin{equation}
  \Depth(\mathbf{x}, t) : \mathbb{R}^3 \times \mathbb{R} \to \mathbb{R}_{\geq 0}.
\end{equation}
The system minimises $\Depth$ along its trajectory: dynamics is gradient descent
on the partition depth field.

\section{The Partition Lagrangian}
\label{sec:ch5-lagrangian}

\begin{consequence}{Partition Lagrangian}{partitionlagrangian}
  The unique Galilean-invariant, scalar-plus-vector-complete Lagrangian
  coupling a particle of mass parameter $\mu$ to the partition depth field
  $\Depth$ is:
  \begin{equation}\label{eq:partitionlagrangian}
    \Lagr = \tfrac{1}{2}\mu|\dot{\mathbf{x}}|^2
            + \mu\dot{\mathbf{x}} \cdot \Apartition
            - \Depth(\mathbf{x}, t),
  \end{equation}
  where $\Apartition(\mathbf{x}, t)$ is the \emph{partition vector potential},
  defined as the curl-free part of the velocity field that minimises $\Depth$.
\end{consequence}

\begin{proof}[Proof sketch]
  By the completeness theorem for Lagrangians (Helmholtz conditions), the
  most general Lagrangian linear in $\dot{\mathbf{x}}$ and $\Depth$ with
  Galilean symmetry is of the form~\eqref{eq:partitionlagrangian}.
  Galilean invariance forbids terms of order $|\dot{\mathbf{x}}|^0$ other
  than $\Depth$ and terms of order $|\dot{\mathbf{x}}|^2$ other than the
  kinetic term.
\end{proof}

\section{The Partition Equation of Motion}
\label{sec:ch5-eom}

Applying the Euler--Lagrange equations to~\eqref{eq:partitionlagrangian}:
\begin{equation}\label{eq:lorentz-shape}
  \mu\ddot{\mathbf{x}} = \mu\dot{\mathbf{x}} \times
  (\nabla \times \Apartition) - \nabla\Depth,
\end{equation}
where $\mathbf{B}_\mathcal{M} = \nabla \times \Apartition$ plays the role of
the magnetic field and $-\nabla\Depth$ plays the role of the electric force.

Crucially, this equation has the \emph{shape} of the Lorentz force law but is
\emph{force-free} in the electromagnetic sense: $\Apartition$ and $\Depth$ are
not external fields but geometric properties of the partition structure.
The ``force'' is gradient descent on the partition depth of the bounded phase
space.

\begin{remark}[Force-free ion trajectories]
  In the mass spectrometric context (Part~VIII), this means that an ion
  in an instrument does not experience forces in the traditional sense.
  It follows the path of steepest descent on the partition depth landscape
  created by the instrument's field configuration.
  The ion trajectory is geometrically determined.
  The GPU implementation (Chapter~\ref{ch:gpu}) uses exactly this gradient
  descent formulation, making the shader pipeline a physical observable of
  the partition depth field.
\end{remark}

\bigskip
\begin{tcolorbox}[colback=crownblue!5, colframe=crownblue!60!black,
                  title={\textbf{Chapter 5 Summary}}]
\begin{itemize}[noitemsep]
  \item The partition depth field $\Depth(\mathbf{x},t)$ is a scalar field
    derived from the partition algebra.
  \item The unique Galilean-invariant Lagrangian coupling a particle to
    $\Depth$ has the form $\Lagr = \half\mu|\dot{\mathbf{x}}|^2
    + \mu\dot{\mathbf{x}}\cdot\Apartition - \Depth$.
  \item The Euler--Lagrange equations have the Lorentz-force shape but are
    force-free: dynamics is gradient descent on $\Depth$.
  \item Newton's second law is a special case of this equation when the
    partition depth gradient is the Newtonian potential.
\end{itemize}
\end{tcolorbox}
